\documentclass[12pt,a4paper]{article}

\usepackage[utf8]{inputenc}
\usepackage[T1]{fontenc}
\usepackage[spanish]{babel}
\usepackage{geometry}
\usepackage{graphicx}
\usepackage{listings}
\usepackage{xcolor}
\usepackage{titlesec}
\usepackage{fancyhdr}
\usepackage{hyperref}
\usepackage{tikz}
\usetikzlibrary{positioning, arrows.meta, shapes.geometric}
\usepackage{amsmath}
\usepackage{float}

\geometry{left=2.5cm, right=2.5cm, top=3cm, bottom=3cm}

\definecolor{codeblue}{rgb}{0.0, 0.35, 0.65}
\definecolor{codegreen}{rgb}{0.0, 0.5, 0.2}
\definecolor{codegray}{rgb}{0.5, 0.5, 0.5}
\definecolor{codepurple}{rgb}{0.52, 0.0, 0.52}
\definecolor{backcolour}{rgb}{0.97, 0.97, 0.97}

\lstdefinestyle{pythonstyle}{
    backgroundcolor=\color{backcolour},
    commentstyle=\color{codegreen}\itshape,
    keywordstyle=\color{codeblue}\bfseries,
    stringstyle=\color{codepurple},
    numberstyle=\tiny\color{codegray},
    basicstyle=\ttfamily\footnotesize,
    breakatwhitespace=false,
    breaklines=true,
    captionpos=b,
    keepspaces=true,
    numbers=left,
    numbersep=5pt,
    showspaces=false,
    showstringspaces=false,
    showtabs=false,
    tabsize=4,
    language=Python,
    frame=single,
    rulecolor=\color{codeblue},
    morekeywords={True, False, None, self}
}

\pagestyle{fancy}
\fancyhf{}
\rhead{\small POO 2026-2S -- Actividad 1}
\lhead{\small Universidad Nacional de Colombia}
\cfoot{\thepage}

\titleformat{\section}[block]{\large\bfseries\color{codeblue}}{}{0em}{}[\titlerule]
\titleformat{\subsection}[block]{\normalsize\bfseries}{}{0em}{}

\hypersetup{
    colorlinks=true,
    linkcolor=codeblue,
    urlcolor=codeblue,
    pdftitle={Actividad 1 - POO 2026-2S},
    pdfauthor={Yuliana Stefan\'ia Guzm\'an Montoya}
}

\begin{document}

\begin{titlepage}
    \centering
    \vspace*{1.5cm}

    {\LARGE \textbf{Universidad Nacional de Colombia}}\\[0.4cm]
    {\large Sede Medell\'in -- Facultad de Minas}\\[0.2cm]
    {\large Departamento de Ciencias de la Computaci\'on y la Decisi\'on}\\[2cm]

    \rule{\linewidth}{1.5pt}\\[0.4cm]
    {\huge \textbf{Actividad 1}}\\[0.3cm]
    {\Large Programaci\'on Orientada a Objetos}\\[0.3cm]
    \rule{\linewidth}{1.5pt}\\[2cm]

    \begin{tabular}{rl}
        \textbf{Asignatura:}  & Programaci\'on Orientada a Objetos \\[0.3cm]
        \textbf{C\'odigo:}    & 2026-2S \\[0.3cm]
        \textbf{Actividad:}   & Actividad 1 -- Individual (10\%) \\[0.3cm]
        \textbf{Estudiante:}  & Yuliana Stefan\'ia Guzm\'an Montoya \\[0.3cm]
        \textbf{Docente:}     & Walter Hugo Arboleda Mazo \\[0.3cm]
        \textbf{Fecha:}       & 17 de septiembre de 2026 \\[0.3cm]
    \end{tabular}

    \vfill

    {\large \textbf{Referencia:}}\\[0.3cm]
    Libro de L\'ogica de Programaci\'on de Efra\'in Oviedo.

    \vspace{1cm}
    \textbf{Repositorio GitHub:}\\
    \url{https://github.com/yguzmanm-sys/POO-2026-2S}
\end{titlepage}

\tableofcontents
\newpage

% ===============================================================
\section{Ejercicio Resuelto No~4}

\subsection{Enunciado}
A la mam\'a de Juan le preguntan su edad, y contesta: tengo 3 hijos, preg\'untele a Juan su edad. Alberto tiene 2/3 de la edad de Juan, Ana tiene 4/3 de la edad de Juan y mi edad es la suma de las tres. Hacer un algoritmo que muestre la edad de los cuatro.

\subsection{C\'odigo fuente -- \texttt{ejercicio\_resuelto4\_edades.py}}
\noindent \textbf{Enlace al c\'odigo en GitHub:}\\
\url{https://github.com/yguzmanm-sys/POO-2026-2S/blob/main/actividad%201/ejercicio_resuelto4_edades.py}\\[0.3cm]

\begin{lstlisting}[style=pythonstyle, caption={Familia de Juan -- ejercicio\_resuelto4\_edades.py}]
class FamiliaJuan:
    def __init__(self, edad_juan: float):
        self.edad_juan = edad_juan

    def calcular_edades(self):
        alberto = (2 / 3) * self.edad_juan
        ana = (4 / 3) * self.edad_juan
        mama = self.edad_juan + alberto + ana
        return self.edad_juan, alberto, ana, mama

    def imprimir_edades(self):
        juan, alberto, ana, mama = self.calcular_edades()
        print(f"Edad de Juan: {juan:.0f}")
        print(f"Edad de Alberto: {alberto:.0f}")
        print(f"Edad de Ana: {ana:.0f}")
        print(f"Edad de la mama: {mama:.0f}")

if __name__ == "__main__":
    familia = FamiliaJuan(9)
    familia.imprimir_edades()
\end{lstlisting}

\subsection{Diagrama de clases}

\begin{figure}[H]
\centering
\begin{tikzpicture}[
    class/.style={draw=codeblue, fill=backcolour, thick, minimum width=6cm, text centered, font=\ttfamily\small},
    section/.style={draw=codeblue, fill=blue!8, thick, minimum width=6cm, text centered, font=\ttfamily\small},
]
\node[section, minimum height=0.8cm] (title) {\textbf{FamiliaJuan}};
\node[class, minimum height=0.9cm, below=0pt of title] (attrs) {
    \begin{tabular}{l}
    -- edad\_juan : float \\
    \end{tabular}
};
\node[class, minimum height=1.6cm, below=0pt of attrs] (methods) {
    \begin{tabular}{l}
    + \_\_init\_\_(edad\_juan:float) \\
    + calcular\_edades() : tuple \\
    + imprimir\_edades() \\
    \end{tabular}
};
\end{tikzpicture}
\end{figure}

\newpage
% ===============================================================
\section{Ejercicio Resuelto No~5}

\subsection{Enunciado}
Hacer un seguimiento (prueba de escritorio) del siguiente grupo de instrucciones.
\begin{verbatim}
INICIO
  SUMA = 0
  X = 20
  SUMA = SUMA + X
  Y = 40
  X = X + Y ** 2
  SUMA = SUMA + X / Y
  ESCRIBA: "EL VALOR DE LA SUMA ES:", SUMA
FIN_INICIO
\end{verbatim}

\subsection{C\'odigo fuente -- \texttt{ejercicio\_resuelto5\_seguimiento.py}}
\noindent \textbf{Enlace al c\'odigo en GitHub:}\\
\url{https://github.com/yguzmanm-sys/POO-2026-2S/blob/main/actividad%201/ejercicio_resuelto5_seguimiento.py}\\[0.3cm]

\begin{lstlisting}[style=pythonstyle, caption={Seguimiento -- ejercicio\_resuelto5\_seguimiento.py}]
class SeguimientoInstrucciones:
    def __init__(self):
        self.suma = 0
        self.x = 0
        self.y = 0

    def ejecutar_instrucciones(self):
        self.suma = 0
        self.x = 20
        self.suma = self.suma + self.x
        self.y = 40
        self.x = self.x + self.y ** 2
        self.suma = self.suma + self.x / self.y
        return self.suma

    def imprimir_resultado(self):
        resultado = self.ejecutar_instrucciones()
        print(f"EL VALOR DE LA SUMA ES: {resultado}")

if __name__ == "__main__":
    seguimiento = SeguimientoInstrucciones()
    seguimiento.imprimir_resultado()
\end{lstlisting}

\subsection{Diagrama de clases}

\begin{figure}[H]
\centering
\begin{tikzpicture}[
    class/.style={draw=codeblue, fill=backcolour, thick, minimum width=6cm, text centered, font=\ttfamily\small},
    section/.style={draw=codeblue, fill=blue!8, thick, minimum width=6cm, text centered, font=\ttfamily\small},
]
\node[section, minimum height=0.8cm] (title) {\textbf{SeguimientoInstrucciones}};
\node[class, minimum height=1.2cm, below=0pt of title] (attrs) {
    \begin{tabular}{l}
    -- suma : float \\
    -- x : float \\
    -- y : float \\
    \end{tabular}
};
\node[class, minimum height=1.6cm, below=0pt of attrs] (methods) {
    \begin{tabular}{l}
    + \_\_init\_\_() \\
    + ejecutar\_instrucciones() : float \\
    + imprimir\_resultado() \\
    \end{tabular}
};
\end{tikzpicture}
\end{figure}

\newpage
% ===============================================================
\section{Ejercicio Propuesto No~12}

\subsection{Enunciado}
Un empleado trabaja 48 horas en la semana a raz\'on de \$5.000 hora. El porcentaje de retenci\'on en la fuente es del 12,5\% del salario bruto. Se desea saber cu\'al es el salario bruto, la retenci\'on en la fuente y el salario neto del trabajador.

\subsection{C\'odigo fuente -- \texttt{ejercicio\_propuesto12\_salario.py}}
\noindent \textbf{Enlace al c\'odigo en GitHub:}\\
\url{https://github.com/yguzmanm-sys/POO-2026-2S/blob/main/actividad%201/ejercicio_propuesto12_salario.py}\\[0.3cm]

\begin{lstlisting}[style=pythonstyle, caption={Salario Empleado -- ejercicio\_propuesto12\_salario.py}]
class Empleado:
    def __init__(self, horas_semana: float, valor_hora: float, retencion_fuente_porcentaje: float):
        self.horas_semana = horas_semana
        self.valor_hora = valor_hora
        self.retencion_fuente_porcentaje = retencion_fuente_porcentaje

    def calcular_salario_bruto(self):
        return self.horas_semana * self.valor_hora

    def calcular_retencion_fuente(self):
        return self.calcular_salario_bruto() * (self.retencion_fuente_porcentaje / 100)

    def calcular_salario_neto(self):
        return self.calcular_salario_bruto() - self.calcular_retencion_fuente()

    def imprimir_informacion(self):
        bruto = self.calcular_salario_bruto()
        retencion = self.calcular_retencion_fuente()
        neto = self.calcular_salario_neto()
        print(f"Salario Bruto: ${bruto:,.2f}")
        print(f"Retencion en la fuente: ${retencion:,.2f}")
        print(f"Salario Neto: ${neto:,.2f}")

if __name__ == "__main__":
    empleado = Empleado(48, 5000, 12.5)
    empleado.imprimir_informacion()
\end{lstlisting}

\subsection{Diagrama de clases}

\begin{figure}[H]
\centering
\begin{tikzpicture}[
    class/.style={draw=codeblue, fill=backcolour, thick, minimum width=8cm, text centered, font=\ttfamily\small},
    section/.style={draw=codeblue, fill=blue!8, thick, minimum width=8cm, text centered, font=\ttfamily\small},
]
\node[section, minimum height=0.8cm] (title) {\textbf{Empleado}};
\node[class, minimum height=1.2cm, below=0pt of title] (attrs) {
    \begin{tabular}{l}
    -- horas\_semana : float \\
    -- valor\_hora : float \\
    -- retencion\_fuente\_porcentaje : float \\
    \end{tabular}
};
\node[class, minimum height=2.0cm, below=0pt of attrs] (methods) {
    \begin{tabular}{l}
    + \_\_init\_\_(horas:float, valor:float, retencion:float) \\
    + calcular\_salario\_bruto() : float \\
    + calcular\_retencion\_fuente() : float \\
    + calcular\_salario\_neto() : float \\
    + imprimir\_informacion() \\
    \end{tabular}
};
\end{tikzpicture}
\end{figure}

\newpage
% ===============================================================
\section{Ejercicio Propuesto No~14}

\subsection{Enunciado}
Elabore un algoritmo que lea un n\'umero y obtenga su cuadrado y su cubo.

\subsection{C\'odigo fuente -- \texttt{ejercicio\_propuesto14\_potencias.py}}
\noindent \textbf{Enlace al c\'odigo en GitHub:}\\
\url{https://github.com/yguzmanm-sys/POO-2026-2S/blob/main/actividad%201/ejercicio_propuesto14_potencias.py}\\[0.3cm]

\begin{lstlisting}[style=pythonstyle, caption={Cuadrado y Cubo -- ejercicio\_propuesto14\_potencias.py}]
class Numero:
    def __init__(self, valor: float):
        self.valor = valor

    def obtener_cuadrado(self):
        return self.valor ** 2

    def obtener_cubo(self):
        return self.valor ** 3

    def imprimir_resultados(self):
        print(f"Numero: {self.valor}")
        print(f"Cuadrado: {self.obtener_cuadrado()}")
        print(f"Cubo: {self.obtener_cubo()}")

if __name__ == "__main__":
    numero = Numero(5)
    numero.imprimir_resultados()
\end{lstlisting}

\subsection{Diagrama de clases}

\begin{figure}[H]
\centering
\begin{tikzpicture}[
    class/.style={draw=codeblue, fill=backcolour, thick, minimum width=6cm, text centered, font=\ttfamily\small},
    section/.style={draw=codeblue, fill=blue!8, thick, minimum width=6cm, text centered, font=\ttfamily\small},
]
\node[section, minimum height=0.8cm] (title) {\textbf{Numero}};
\node[class, minimum height=0.9cm, below=0pt of title] (attrs) {
    \begin{tabular}{l}
    -- valor : float \\
    \end{tabular}
};
\node[class, minimum height=1.6cm, below=0pt of attrs] (methods) {
    \begin{tabular}{l}
    + \_\_init\_\_(valor:float) \\
    + obtener\_cuadrado() : float \\
    + obtener\_cubo() : float \\
    + imprimir\_resultados() \\
    \end{tabular}
};
\end{tikzpicture}
\end{figure}

\newpage
% ===============================================================
\section{Ejercicio Propuesto No~17}

\subsection{Enunciado}
Dado el radio de un c\'irculo. Haga un algoritmo que obtenga el \'area del c\'irculo y la longitud de la circunferencia.

\subsection{C\'odigo fuente -- \texttt{ejercicio\_propuesto17\_circulo.py}}
\noindent \textbf{Enlace al c\'odigo en GitHub:}\\
\url{https://github.com/yguzmanm-sys/POO-2026-2S/blob/main/actividad%201/ejercicio_propuesto17_circulo.py}\\[0.3cm]

\begin{lstlisting}[style=pythonstyle, caption={C\'irculo -- ejercicio\_propuesto17\_circulo.py}]
import math

class Circulo:
    def __init__(self, radio: float):
        self.radio = radio

    def calcular_area(self):
        return math.pi * self.radio ** 2

    def calcular_circunferencia(self):
        return 2 * math.pi * self.radio

    def imprimir_resultados(self):
        print(f"Radio del circulo: {self.radio}")
        print(f"Area del circulo: {self.calcular_area():.4f}")
        print(f"Longitud de la circunferencia: {self.calcular_circunferencia():.4f}")

if __name__ == "__main__":
    circulo = Circulo(10)
    circulo.imprimir_resultados()
\end{lstlisting}

\subsection{Diagrama de clases}

\begin{figure}[H]
\centering
\begin{tikzpicture}[
    class/.style={draw=codeblue, fill=backcolour, thick, minimum width=6.5cm, text centered, font=\ttfamily\small},
    section/.style={draw=codeblue, fill=blue!8, thick, minimum width=6.5cm, text centered, font=\ttfamily\small},
]
\node[section, minimum height=0.8cm] (title) {\textbf{Circulo}};
\node[class, minimum height=0.9cm, below=0pt of title] (attrs) {
    \begin{tabular}{l}
    -- radio : float \\
    \end{tabular}
};
\node[class, minimum height=1.6cm, below=0pt of attrs] (methods) {
    \begin{tabular}{l}
    + \_\_init\_\_(radio:float) \\
    + calcular\_area() : float \\
    + calcular\_circunferencia() : float \\
    + imprimir\_resultados() \\
    \end{tabular}
};
\end{tikzpicture}
\end{figure}

% ---------------------------------------------------------------
\newpage
\section{Repositorio GitHub}

El c\'odigo fuente y la documentaci\'on de la Actividad 1 se encuentran alojados en GitHub:

\begin{itemize}
    \setlength\itemsep{0.5em}
    \item \textbf{Repositorio General:} \\
    \url{https://github.com/yguzmanm-sys/POO-2026-2S}
    \item \textbf{Directorio de la Actividad 1:} \\
    \url{https://github.com/yguzmanm-sys/POO-2026-2S/tree/main/actividad%201}
    \item \textbf{Ejercicio Resuelto No~4 (\texttt{ejercicio\_resuelto4\_edades.py}):} \\
    \url{https://github.com/yguzmanm-sys/POO-2026-2S/blob/main/actividad%201/ejercicio_resuelto4_edades.py}
    \item \textbf{Ejercicio Resuelto No~5 (\texttt{ejercicio\_resuelto5\_seguimiento.py}):} \\
    \url{https://github.com/yguzmanm-sys/POO-2026-2S/blob/main/actividad%201/ejercicio_resuelto5_seguimiento.py}
    \item \textbf{Ejercicio Propuesto No~12 (\texttt{ejercicio\_propuesto12\_salario.py}):} \\
    \url{https://github.com/yguzmanm-sys/POO-2026-2S/blob/main/actividad%201/ejercicio_propuesto12_salario.py}
    \item \textbf{Ejercicio Propuesto No~14 (\texttt{ejercicio\_propuesto14\_potencias.py}):} \\
    \url{https://github.com/yguzmanm-sys/POO-2026-2S/blob/main/actividad%201/ejercicio_propuesto14_potencias.py}
    \item \textbf{Ejercicio Propuesto No~17 (\texttt{ejercicio\_propuesto17\_circulo.py}):} \\
    \url{https://github.com/yguzmanm-sys/POO-2026-2S/blob/main/actividad%201/ejercicio_propuesto17_circulo.py}
\end{itemize}

\end{document}
